\chapter{Finite Calculus}

\begin{center}\itshape
Everyone knows $1+2+\dots+n=\frac{n(n+1)}2$. What is $1^5+2^5+\dots+n^5$,\\
and is there a method that would have found it, rather than a formula to remember?
\end{center}

In his \emph{Ars Conjectandi}, published in 1713, Jacob Bernoulli could not resist a boast. With his new method, he wrote, it had taken him
``less than half of a quarter of an hour'' to find that the tenth powers of the first thousand integers add up to
\[
91\,409\,924\,241\,424\,243\,424\,241\,924\,242\,500.
\]
(A computer confirms every digit.) His secret was not speed with arithmetic but a calculus for sums, as complete as Newton's calculus for areas.
This chapter builds it.

Calculus computes $\int_a^bf$ by finding a function $F$ with $F'=f$. Sums have the same theory, with the
\emph{difference} $\Delta F(x)=F(x+1)-F(x)$ in place of the derivative: to compute $\sum_{k=a}^{b-1}f(k)$, find $F$
with $\Delta F=f$, and the sum is $F(b)-F(a)$. Everything else follows the continuous model, with small twists.
Powers $x^m$ are replaced by \emph{falling powers} $x(x-1)\cdots(x-m+1)$, which behave perfectly under $\Delta$;
integration by parts becomes \emph{summation by parts}; the exponential $e^x$ becomes $2^x$ and the logarithm
becomes the harmonic number $H_x$.

The chapter develops this \emph{finite calculus}, then uses it on the oldest problem of its kind, the sums of
powers $\sum k^p$. Two families of numbers do the translation between ordinary and falling powers: the
\emph{Stirling numbers}, which count partitions of sets and cycles of permutations, and the \emph{Bernoulli
numbers}, which give Faulhaber's formula for $\sum k^p$ and will return in Chapter~3 with $\zeta(2k)$ and in
Chapter~13 with the Euler--Maclaurin formula. The chapter ends with Newton's interpolation formula, the principle
of Babbage's difference engine.

\begin{think}
The exponential is the function equal to its own derivative. Which function $f$ on the integers satisfies
$\Delta f=f$, with $f(0)=1$?
\end{think}

\section{The difference operator}

\begin{definition}
For a function $f$ defined on $\R$ (or on $\Z$), its \emph{difference} is $\Delta f(x)=f(x+1)-f(x)$. For integers
$a\le b$ we write $\sum_a^bf=\sum_{k=a}^{b-1}f(k)$, a sum over the half-open range $a\le k<b$ (it is $0$ when $a=b$).
\end{definition}

The half-open convention makes ranges add up, $\sum_a^bf+\sum_b^cf=\sum_a^cf$, exactly as for integrals, and it
makes the next theorem exact.

\begin{theorem}[fundamental theorem of finite calculus]\label{thm:ftfc}
If $\Delta F=f$ on the integers, then $\sum_a^bf=F(b)-F(a)$ for all integers $a\le b$. Conversely, if
$\Delta F=\Delta G$ on $\Z$, then $F-G$ is constant on $\Z$.
\end{theorem}
\begin{proof}
$\sum_{k=a}^{b-1}\bigl(F(k+1)-F(k)\bigr)$ telescopes to $F(b)-F(a)$. If $\Delta(F-G)=0$ then $(F-G)(k+1)=(F-G)(k)$ for
every integer $k$, so $F-G$ takes the same value at all integers.
\end{proof}

So every sum is a problem of finding an \emph{antidifference}. The first family to try is the powers, and they
fail: $\Delta x^2=2x+1$, $\Delta x^3=3x^2+3x+1$, not as clean as $(x^m)'=mx^{m-1}$. The remedy is to change the powers.

\begin{definition}
The \emph{falling powers} are $x^{\underline 0}=1$ and $x^{\underline m}=x(x-1)\cdots(x-m+1)$ for $m\ge1$ ($m$ factors).
The \emph{negative falling powers} are $x^{\underline{-m}}=\dfrac1{(x+1)(x+2)\cdots(x+m)}$ for $m\ge1$.
\end{definition}

For an integer $n\ge0$, $n^{\underline m}=\frac{n!}{(n-m)!}$ counts the ordered choices of $m$ objects among $n$ when
$m\le n$, and $\binom xm=\frac{x^{\underline m}}{m!}$ is a polynomial in $x$ of degree $m$. The negative powers are
designed so that $x^{\underline{m+n}}=x^{\underline m}(x-m)^{\underline n}$ holds for all integers $m$, $n$, as one
checks by cancelling factors.

\begin{proposition}\label{prop:falling}
For every integer $m$ (positive, zero or negative), $\Delta x^{\underline m}=m\,x^{\underline{m-1}}$. In particular,
for $m\ge0$ and integers $n\ge0$,
\[
\sum_{k=0}^{n-1}k^{\underline m}=\frac{n^{\underline{m+1}}}{m+1},\qq\sum_{k=0}^{n-1}\binom km=\binom n{m+1}.
\]
\end{proposition}
\begin{proof}
For $m\ge1$, $x^{\underline m}=x^{\underline{m-1}}(x-m+1)$ and $(x+1)^{\underline m}=(x+1)x^{\underline{m-1}}$, so
$\Delta x^{\underline m}=x^{\underline{m-1}}\bigl((x+1)-(x-m+1)\bigr)=m\,x^{\underline{m-1}}$. For $m=0$ both sides are
$0$. For $m=-r$ with $r\ge1$, put $P(x)=(x+2)\cdots(x+r)$, so that $x^{\underline{-r}}=\frac1{(x+1)P(x)}$ and
$(x+1)^{\underline{-r}}=\frac1{P(x)(x+r+1)}$. Then
\[
\begin{aligned}\Delta x^{\underline{-r}}&=\frac{(x+1)-(x+r+1)}{(x+1)P(x)(x+r+1)}\eqbr=\frac{-r}{(x+1)\cdots(x+r+1)}=-r\,x^{\underline{-r-1}}.\end{aligned}
\]
The sums follow from Theorem~\ref{thm:ftfc} with $F(x)=\frac{x^{\underline{m+1}}}{m+1}$, which vanishes at $0$; the
second is the first divided by $m!$.
\end{proof}

The second sum is the \emph{hockey-stick identity} of Pascal's triangle. Negative powers telescope the sums of
reciprocals of products: with $m=-2$, $\Delta x^{\underline{-1}}=-x^{\underline{-2}}$, so
\[
\sum_{k=0}^{n-1}\frac1{(k+1)(k+2)}=\Bigl[-x^{\underline{-1}}\Bigr]_0^n=1-\frac1{n+1}.
\]
To reach $x^{\underline{-1}}$ one would need $m=0$, and the rule then gives $\Delta x^{\underline 0}=0$, which is true but useless: $x^{\underline{-1}}=\frac1{x+1}$
is the one power that is not a difference of a power, just as $\frac1x$ is the one power that is not a derivative of a
power. Its antidifference is the \emph{harmonic number} $H_x$: for integers, $H_n=1+\frac12+\dots+\frac1n$ satisfies
$\Delta H_n=\frac1{n+1}=n^{\underline{-1}}$, so $H$ plays the part of the logarithm. And $\Delta2^x=2^x$: the function
$2^x$ is the discrete exponential that answers the question opening the chapter, and more generally
$\Delta c^x=(c-1)c^x$. Table~\ref{tab:dictionary} collects the dictionary.

\begin{table}[ht]
\centering
\begin{tabular}{ll}
\hline
continuous & discrete\\
\hline
$f'(x)$ & $\Delta f(x)=f(x+1)-f(x)$\\
$\int_a^bf(x)\dd x$ & $\sum_a^bf=\sum_{a\le k<b}f(k)$\\
$(x^m)'=mx^{m-1}$ & $\Delta x^{\underline m}=mx^{\underline{m-1}}$\\
$(\ln x)'=x^{-1}$ & $\Delta H_x=x^{\underline{-1}}$\\
$(e^x)'=e^x$ & $\Delta2^x=2^x$\\
$(uv)'=u'v+uv'$ & $\Delta(uv)=u\,\Delta v+(\Delta u)\,Ev$\\
\hline
\end{tabular}
\caption{Calculus and finite calculus. $Ev(x)=v(x+1)$ is the shift.}
\label{tab:dictionary}
\end{table}

\begin{figure}[ht]
\centering
\fig{ch01-staircase}
\caption{The sum $\sum_{k=0}^{n-1}k^2$ (the bars, here $n=8$) against the integral $\int_0^nx^2\dd x=\frac{n^3}3$ (the area
under the curve). The sum is $\frac{n^{\underline3}}3+\frac{n^{\underline2}}2=\frac{n^3}3-\frac{n^2}2+\frac n6$: the falling powers
give it exactly, and the integral only gives its leading term.}
\label{fig:staircase}
\end{figure}

\section{Summation by parts}

The product rule has a discrete form with a shift in it: if $Ev(x)=v(x+1)$, then
\[
\begin{aligned}\Delta(uv)(x)&=u(x+1)v(x+1)-u(x)v(x)\eqbr=u(x)\,\Delta v(x)+\Delta u(x)\,v(x+1),\end{aligned}
\]
since the right-hand side is $u(x)v(x+1)-u(x)v(x)+u(x+1)v(x+1)-u(x)v(x+1)$. Summing over $a\le k<b$ with
Theorem~\ref{thm:ftfc} gives the discrete integration by parts.

\begin{proposition}[summation by parts]\label{prop:byparts}
For functions $u$, $v$ on the integers and integers $a\le b$,
\[
\sum_a^bu\,\Delta v=\bigl[uv\bigr]_a^b-\sum_a^b(\Delta u)\,Ev,
\]
where $[uv]_a^b=u(b)v(b)-u(a)v(a)$.
\end{proposition}

\begin{example}\label{ex:k2k}
To sum $\sum_{k=0}^{n-1}k2^k$, take $u(k)=k$ and $\Delta v=2^k$, so $v(k)=2^k$, $\Delta u=1$ and $Ev(k)=2^{k+1}$:
\[
\begin{aligned}\sum_{k=0}^{n-1}k2^k&=n2^n-\sum_{k=0}^{n-1}2^{k+1}\eqbr=n2^n-2(2^n-1)\eqbr=(n-2)2^n+2.\end{aligned}
\]
For $n=3$: $0+2+8=10=1\cdot8+2$. With $u=H_k$ and $\Delta v=1$, so $v(k)=k$, the same rule gives
\[
\sum_{k=0}^{n-1}H_k=nH_n-\sum_{k=0}^{n-1}\frac{k+1}{k+1}=nH_n-n.
\]
\end{example}

Abel's form of summation by parts is the one used for series. Write $A_k=a_1+\dots+a_k$ and $A_0=0$, so that
$a_k=\Delta A_{k-1}$.

\begin{proposition}[Abel summation]\label{prop:abel}
For numbers $a_k$, $b_k$ and $n\ge1$,
\[
\sum_{k=1}^na_kb_k=A_nb_n-\sum_{k=1}^{n-1}A_k(b_{k+1}-b_k).
\]
\end{proposition}
\begin{proof}
$\sum_{k=1}^na_kb_k=\sum_{k=1}^n(A_k-A_{k-1})b_k=\sum_{k=1}^nA_kb_k-\sum_{k=0}^{n-1}A_kb_{k+1}$. Since $A_0=0$, this
is $A_nb_n+\sum_{k=1}^{n-1}A_k(b_k-b_{k+1})$.
\end{proof}

\begin{theorem}[Dirichlet's test]\label{thm:dirichlet}
Let $(a_k)$ be complex numbers whose partial sums $A_n$ are bounded, $|A_n|\le M$, and let $(b_k)$ be real numbers
that decrease to $0$. Then $\sum a_kb_k$ converges, and $\bigl|\sum_{k\ge1}a_kb_k\bigr|\le Mb_1$.
\end{theorem}
\begin{proof}
In Proposition~\ref{prop:abel}, $|A_nb_n|\le Mb_n\to0$. The series $\sum_kA_k(b_{k+1}-b_k)$ converges absolutely,
since $|A_k(b_{k+1}-b_k)|\le M(b_k-b_{k+1})$ and $\sum_{k=1}^{n-1}(b_k-b_{k+1})=b_1-b_n\le b_1$. So the partial sums
$\sum_{k=1}^na_kb_k$ converge, and the same bounds give $\bigl|\sum_{k=1}^na_kb_k\bigr|\le Mb_n+M(b_1-b_n)=Mb_1$.
\end{proof}

With $a_k=(-1)^{k+1}$ this is the alternating series test of Volume~1. The more interesting case is trigonometric.

\begin{proposition}\label{prop:sinsum}
For $x\notin2\pi\Z$ and $n\ge1$,
\[
\sum_{k=1}^n\sin kx=\frac{\sin\frac{nx}2\,\sin\frac{(n+1)x}2}{\sin\frac x2}.
\]
In particular $\bigl|\sum_{k=1}^n\sin kx\bigr|\le\frac1{|\sin(x/2)|}$, and $\sum_{k\ge1}\frac{\sin kx}k$ converges for every real $x$.
\end{proposition}
\begin{proof}
The identity $2\sin\frac x2\sin kx=\cos\bigl(k-\frac12\bigr)x-\cos\bigl(k+\frac12\bigr)x$ makes the sum telescope:
\[
\begin{aligned}2\sin\frac x2\sum_{k=1}^n\sin kx&=\cos\frac x2-\cos\Bigl(n+\frac12\Bigr)x\eqbr=2\sin\frac{nx}2\sin\frac{(n+1)x}2,\end{aligned}
\]
the last step by $\cos p-\cos q=2\sin\frac{q+p}2\sin\frac{q-p}2$. For $x\notin2\pi\Z$, Theorem~\ref{thm:dirichlet} with
$a_k=\sin kx$ and $b_k=\frac1k$ gives the convergence; for $x\in2\pi\Z$ every term is $0$.
\end{proof}

Chapter~8 shows that the sum of this series is $\frac{\pi-x}2$ for $0<x<2\pi$. The convergence is not absolute: the
terms are not summable in absolute value, and only the cancellations that summation by parts exposes make the series
converge.

\section{Stirling numbers}

Falling powers sum easily and ordinary powers do not, so we need to convert between the two. The coefficients are
classical combinatorial numbers.

\begin{definition}
For integers $n,k\ge0$, the \emph{Stirling number of the second kind} $\stirlingii nk$ is the number of partitions of an
$n$-element set into $k$ non-empty blocks. The \emph{(unsigned) Stirling number of the first kind} $\stirlingi nk$ is the
number of permutations of an $n$-element set with exactly $k$ cycles. By convention $\stirlingii00=\stirlingi00=1$, and
$\stirlingii n0=\stirlingi n0=0$ for $n\ge1$.
\end{definition}

For example $\stirlingii42=7$: the four partitions $\{a\}\{bcd\}$, $\{b\}\{acd\}$, $\{c\}\{abd\}$, $\{d\}\{abc\}$ and
the three $\{ab\}\{cd\}$, $\{ac\}\{bd\}$, $\{ad\}\{bc\}$. And $\stirlingi42=11$: eight permutations of type
$3+1$ and three of type $2+2$. Volume~4, Chapter~2, computed the expected number of cycles of a random permutation,
which is $\frac1{n!}\sum_kk\stirlingi nk=H_n$.

\begin{proposition}\label{prop:stirlingrec}
For $n\ge0$ and $k\ge1$,
\[
\begin{gathered}\stirlingii{n+1}k=k\stirlingii nk+\stirlingii n{k-1},\qq\stirlingi{n+1}k=n\stirlingi nk+\stirlingi n{k-1}.\end{gathered}
\]
\end{proposition}
\begin{proof}
Look at the element $n+1$. In a partition of $\{1,\dots,n+1\}$ into $k$ blocks, either $\{n+1\}$ is a block, and the
rest is a partition of $\{1,\dots,n\}$ into $k-1$ blocks, or $n+1$ was added to one of the $k$ blocks of a partition
of $\{1,\dots,n\}$ into $k$ blocks. In a permutation of $\{1,\dots,n+1\}$ with $k$ cycles, either $n+1$ is a fixed
point, and the rest is a permutation of $\{1,\dots,n\}$ with $k-1$ cycles, or deleting $n+1$ from its cycle (sending its
predecessor to its successor) gives a permutation of $\{1,\dots,n\}$ with $k$ cycles; conversely $n+1$ can be inserted
after any of the $n$ elements, in $n$ ways.
\end{proof}

\begin{table}[ht]
\centering
\begin{tabular}{c|rrrrrr}
$\stirlingii nk$ & $k=1$ & $2$ & $3$ & $4$ & $5$ & $6$\\
\hline
$n=1$ & 1 &&&&&\\
$2$ & 1 & 1 &&&&\\
$3$ & 1 & 3 & 1 &&&\\
$4$ & 1 & 7 & 6 & 1 &&\\
$5$ & 1 & 15 & 25 & 10 & 1 &\\
$6$ & 1 & 31 & 90 & 65 & 15 & 1\\
\end{tabular}

\medskip
\begin{tabular}{c|rrrrrr}
$\stirlingi nk$ & $k=1$ & $2$ & $3$ & $4$ & $5$ & $6$\\
\hline
$n=1$ & 1 &&&&&\\
$2$ & 1 & 1 &&&&\\
$3$ & 2 & 3 & 1 &&&\\
$4$ & 6 & 11 & 6 & 1 &&\\
$5$ & 24 & 50 & 35 & 10 & 1 &\\
$6$ & 120 & 274 & 225 & 85 & 15 & 1\\
\end{tabular}
\caption{Stirling numbers of the second kind (left) and of the first kind (right).}
\label{tab:stirling}
\end{table}

Write $x^{\overline n}=x(x+1)\cdots(x+n-1)$ for the \emph{rising powers}, with $x^{\overline0}=1$. Note
$x^{\underline n}=(-1)^n(-x)^{\overline n}$.

\begin{theorem}\label{thm:stirling}
For every $n\ge0$, as polynomials in $x$,
\[
\begin{gathered}x^n=\sum_{k=0}^n\stirlingii nk x^{\underline k},\qq x^{\overline n}=\sum_{k=0}^n\stirlingi nk x^k,\\ x^{\underline n}=\sum_{k=0}^n(-1)^{n-k}\stirlingi nk x^k.\end{gathered}
\]
\end{theorem}
\begin{proof}
First identity. Both sides are polynomials, so it is enough to check it at the infinitely many integers $x=m\ge1$
(two polynomials that agree at infinitely many points are equal, Volume~1, Chapter~6). The left side $m^n$ counts the maps
from $\{1,\dots,n\}$ to $\{1,\dots,m\}$. Sort them by the partition of $\{1,\dots,n\}$ into the non-empty fibres of the
map: for a partition into $k$ blocks there are $m^{\underline k}$ maps with exactly these fibres, one for each injective
assignment of distinct values to the blocks. Summing over $k$ gives $\sum_k\stirlingii nk m^{\underline k}$.

Second identity, by induction on $n$; it holds for $n=0$. If it holds for $n$, then
\[
\begin{aligned}x^{\overline{n+1}}&=(x+n)x^{\overline n}=\sum_k\stirlingi nk x^{k+1}+n\sum_k\stirlingi nk x^k\eqbr=\sum_k\Bigl(\stirlingi n{k-1}+n\stirlingi nk\Bigr)x^k,\end{aligned}
\]
which is $\sum_k\stirlingi{n+1}k x^k$ by Proposition~\ref{prop:stirlingrec}. The third identity is the second at $-x$,
multiplied by $(-1)^n$.
\end{proof}

Combined with Proposition~\ref{prop:falling}, the first identity sums any power:
\[
\sum_{k=0}^{n-1}k^p=\sum_{j=0}^p\stirlingii pj\frac{n^{\underline{j+1}}}{j+1}.
\]
For $p=2$ this is $\frac{n^{\underline2}}2+\frac{n^{\underline3}}3=\frac{(n-1)n(2n-1)}6$, and for $p=3$ it is
$\frac{n^{\underline2}}2+n^{\underline3}+\frac{n^{\underline4}}4=\Bigl(\frac{n(n-1)}2\Bigr)^2$.

\begin{think}
$1+2+\dots+(n-1)=\frac{n(n-1)}2$ and $1^3+2^3+\dots+(n-1)^3=\bigl(\frac{n(n-1)}2\bigr)^2$. Is the sum of the fifth powers
also a polynomial in $\frac{n(n-1)}2$?
\end{think}

\section{Bernoulli numbers and sums of powers}

The Stirling route gives each sum of powers, but not a single formula for all of them. Jacob Bernoulli found one in
1713, through the numbers now named after him. Their natural home is a generating function, and we use the language
of formal power series, which Chapter~2 develops further: a formal power series $\sum_{n\ge0}c_nt^n$ is just its
sequence of coefficients, added termwise and multiplied by the Cauchy product
$\bigl(\sum a_nt^n\bigr)\bigl(\sum b_nt^n\bigr)=\sum_n\bigl(\sum_{k=0}^na_kb_{n-k}\bigr)t^n$, without any question of
convergence. A series with $c_0\ne0$ has an inverse: the equations $\sum_{k=0}^nc_kd_{n-k}=[n=0]$ determine $d_0=\frac1{c_0}$,
then $d_1$, $d_2$, \dots one after the other. We write $e^{at}=\sum_n\frac{a^nt^n}{n!}$; the binomial theorem gives
$e^{at}e^{bt}=e^{(a+b)t}$, since the coefficient of $\frac{t^n}{n!}$ on the left is $\sum_k\binom nka^kb^{n-k}$.

\begin{definition}
The \emph{Bernoulli numbers} $B_n$ are defined by the formal identity
\[
\frac t{e^t-1}=\sum_{n\ge0}B_n\frac{t^n}{n!},
\]
where $\frac t{e^t-1}$ is the inverse of the series $\frac{e^t-1}t=\sum_{n\ge0}\frac{t^n}{(n+1)!}$, whose constant term is $1$.
The \emph{Bernoulli polynomials} are defined by $\dfrac{te^{xt}}{e^t-1}=\sum_{n\ge0}B_n(x)\dfrac{t^n}{n!}$, that is,
$B_n(x)=\sum_{k=0}^n\binom nkB_kx^{n-k}$.
\end{definition}

\begin{proposition}\label{prop:bernoulli}
\begin{enumerate}[label=(\alph*)]
\item $B_0=1$, and for $n\ge1$
\[
\sum_{k=0}^{n}\binom{n+1}kB_k=0.
\]
Hence $B_1=-\frac12$, $B_2=\frac16$, $B_4=-\frac1{30}$,
$B_6=\frac1{42}$, $B_8=-\frac1{30}$, $B_{10}=\frac5{66}$, $B_{12}=-\frac{691}{2730}$.
\item $B_n=0$ for every odd $n\ge3$.
\item $B_n(x+1)-B_n(x)=nx^{n-1}$ for $n\ge1$, and $B_n(0)=B_n$.
\end{enumerate}
\end{proposition}
\begin{proof}
(a) Multiply $\sum_kB_k\frac{t^k}{k!}$ by $e^t-1=\sum_{j\ge1}\frac{t^j}{j!}$; the product is $t$. The coefficient of $t^{n+1}$ on the
left is $\sum_{k=0}^n\frac{B_k}{k!\,(n+1-k)!}$, so multiplying by $(n+1)!$ gives $\sum_{k=0}^n\binom{n+1}kB_k=[n=0]$. For $n=0$
this says $B_0=1$; for $n=1$, $1+2B_1=0$; for $n=2$, $1+3B_1+3B_2=0$, so $B_2=\frac16$; and so on.

(b) Let $f(t)=\frac t{e^t-1}+\frac t2=\sum_{n\ne1}B_n\frac{t^n}{n!}$. Its coefficients are those of an even series exactly
when $f(-t)=f(t)$, which is a formal identity to be checked. Multiply both sides by the (non-invertible but harmless)
series $e^t-1$ and use $e^te^{-t}=1$: $(e^t-1)f(t)=t+\frac t2(e^t-1)=\frac t2(e^t+1)$, and
$(e^t-1)f(-t)=(e^t-1)\bigl(\frac{-t}{e^{-t}-1}-\frac t2\bigr)=te^t-\frac t2(e^t-1)=\frac t2(e^t+1)$, using
$\frac{-t}{e^{-t}-1}=\frac{te^t}{e^t-1}$ (multiply numerator and denominator by $-e^t$). So $(e^t-1)\bigl(f(t)-f(-t)\bigr)=0$.
If $g=f(t)-f(-t)$ were non-zero, with lowest term $g_mt^m$, the product $(e^t-1)g$ would have lowest term $g_mt^{m+1}\ne0$.
So $g=0$, and the odd coefficients of $f$, which are the $B_n$ with odd $n\ge3$, vanish.

(c) The difference of the generating functions at $x+1$ and $x$ is $\frac{te^{xt}(e^t-1)}{e^t-1}=te^{xt}=\sum_{n\ge1}nx^{n-1}\frac{t^n}{n!}$.
At $x=0$ the generating function is that of the $B_n$.
\end{proof}

\begin{theorem}[Faulhaber's formula]\label{thm:faulhaber}
For integers $p\ge0$ and $n\ge1$,
\[
\begin{aligned}\sum_{k=0}^{n-1}k^p&=\frac{B_{p+1}(n)-B_{p+1}}{p+1}\eqbr=\frac1{p+1}\sum_{j=0}^p\binom{p+1}jB_jn^{p+1-j}.\end{aligned}
\]
\end{theorem}
\begin{proof}
By Proposition~\ref{prop:bernoulli}(c), $F(x)=\frac{B_{p+1}(x)}{p+1}$ satisfies $\Delta F=x^p$, so Theorem~\ref{thm:ftfc}
gives the first form. Expanding $B_{p+1}(n)=\sum_{j=0}^{p+1}\binom{p+1}jB_jn^{p+1-j}$ and removing the last term
$j=p+1$, which is $B_{p+1}$, gives the second.
\end{proof}

For $p=1$ this is $\frac12(n^2+2B_1n)=\frac{n(n-1)}2$. For $p=5$, with $B_3=B_5=0$,
\[
\sum_{k=0}^{n-1}k^5=\frac{n^6}6-\frac{n^5}2+\frac{5n^4}{12}-\frac{n^2}{12}.
\]
So
\[
\begin{aligned}\sum_{k<1000}k^5&=\frac{10^{18}}6-\frac{10^{15}}2+\frac{5\cdot10^{12}}{12}-\frac{10^6}{12}\eqbr=166\,167\,083\,333\,250\,000\end{aligned}
\]
. Bernoulli was proud of exactly this: he wrote that he had summed the tenth powers of the
first thousand integers ``in half of a quarter of an hour''. The coefficients of Faulhaber's formula are very regular
at the top, $\frac{n^{p+1}}{p+1}-\frac{n^p}2+\frac p{12}n^{p-1}-\cdots$, because $\frac{n^{p+1}}{p+1}$ is the integral and the
correction $-\frac{n^p}2$ is the trapezoid rule; Chapter~13 turns this observation into the Euler--Maclaurin formula.

\begin{remark}
The Bernoulli numbers look tame up to $B_{10}=\frac5{66}$ and then the numerator $691$ appears. Their size is governed
by Euler's formula $\zeta(2k)=\frac{(-1)^{k+1}B_{2k}(2\pi)^{2k}}{2(2k)!}$, proved in Chapter~3: since $\zeta(2k)\to1$,
$|B_{2k}|$ grows like $\frac{2(2k)!}{(2\pi)^{2k}}$, faster than any exponential, so the series $\frac t{e^t-1}$ has radius of
convergence $2\pi$ (its nearest singularities are the poles $\pm2\pi i$, as Chapter~11 explains). Their denominators are
given by the theorem of von Staudt and Clausen (Exercise~\ref{ex:staudt}), and their numerators are linked to Kummer's
work on Fermat's last theorem: the prime $691$ divides the numerator of $B_{12}$, so it is an \emph{irregular} prime.
\end{remark}

\section{Newton's formula and the method of differences}

Iterating $\Delta$ gives $\Delta^nf$. Since $\Delta=E-1$, where $E$ is the shift and $1$ the identity, and $E$ commutes
with itself, the binomial theorem gives
\[
\Delta^nf(x)=\sum_{k=0}^n(-1)^{n-k}\binom nkf(x+k).
\]
(For a careful reader: the operators $E^k$ act on functions and commute, so the expansion of $(E-1)^n$ is the
ordinary algebraic one; or check it by induction on $n$, using Pascal's rule.)

\begin{theorem}[Newton's interpolation formula]\label{thm:newton}
If $f$ is a polynomial of degree at most $d$, then $\Delta^{d+1}f=0$ and, for all $x$,
\[
f(x)=\sum_{k=0}^d\Delta^kf(0)\binom xk.
\]
\end{theorem}
\begin{proof}
By Proposition~\ref{prop:falling}, $\Delta\binom xk=\binom x{k-1}$ for $k\ge1$ and $\Delta\binom x0=0$. The polynomials
$\binom x0,\dots,\binom xd$ have degrees $0,\dots,d$, so they form a basis of the polynomials of degree at most $d$:
write $f=\sum_{k=0}^dc_k\binom xk$. Then $\Delta^jf=\sum_kc_k\binom x{k-j}$ (terms with $k<j$ vanish), whose value at
$0$ is $c_j$, because $\binom0{k-j}=[k=j]$. Also $\Delta^{d+1}f=0$, since each $\Delta$ lowers the degree.
\end{proof}

Newton's formula recovers a polynomial from its values at $0,1,\dots,d$, using only subtractions: the numbers
$\Delta^kf(0)$ are the left edge of the \emph{difference table} (Figure~\ref{fig:difftable}). Read backwards, it computes
a polynomial at $0,1,2,\dots$ using only additions: once the left edge is known, each new row entry is the sum of
two neighbours. Charles Babbage's difference engine, designed in the 1820s to print tables of logarithms and
trigonometric functions free of human error, was a machine for exactly this: it held the differences of a polynomial
approximating the function and added them, column after column.

\begin{figure}[ht]
\centering
\fig{ch01-difftable}
\caption{The difference table of $f(x)=x^3+x+1$. Each entry is the difference of the two above it; the third
differences are constant, equal to $3!=6$. The left edge $1,2,6,6$ gives Newton's formula
$f(x)=1+2\binom x1+6\binom x2+6\binom x3$. The last diagonal is computed from the previous one by additions alone,
from the bottom up: $6$, then $18+6=24$, $38+24=62$, and $f(5)=69+62=131$.}
\label{fig:difftable}
\end{figure}

Applied to $f(x)=x^n$, Newton's formula gives an explicit formula for the Stirling numbers. Comparing
$x^n=\sum_kk!\stirlingii nk\binom xk$ (Theorem~\ref{thm:stirling}) with Theorem~\ref{thm:newton}, by uniqueness of the
coefficients in a basis,
\[
\stirlingii nk=\frac1{k!}\Delta^k\bigl[x^n\bigr]_{x=0}=\frac1{k!}\sum_{j=0}^k(-1)^{k-j}\binom kjj^n.
\]
This is also inclusion--exclusion: $k!\stirlingii nk$ is the number of surjections from an $n$-set onto a $k$-set
(Exercise~\ref{ex:surjections}).

\section{Exercises}

\begin{exercise}[$\star$]\label{ex:k22k}
Show that $\displaystyle\sum_{k=0}^nk^22^k=(n^2-2n+3)2^{n+1}-6$.
\end{exercise}

\begin{exercise}[$\star$]\label{ex:triple}
Compute $\displaystyle\sum_{k=1}^n\frac1{k(k+1)(k+2)}$, and deduce $\displaystyle\sum_{k\ge1}\frac1{k(k+1)(k+2)}=\frac14$.
\end{exercise}

\begin{exercise}[$\star\star$]\label{ex:kHk}
Show that $\displaystyle\sum_{k=0}^{n-1}kH_k=\frac{n(n-1)}2\Bigl(H_n-\frac12\Bigr)$.
\end{exercise}

\begin{exercise}[$\star$]\label{ex:surjections}
(a) Show that the number of surjections from an $n$-set onto a $k$-set is $k!\stirlingii nk$, and, by inclusion--exclusion,
that it is $\sum_{j=0}^k(-1)^{k-j}\binom kjj^n$.
(b) Deduce that $\sum_{j=0}^n(-1)^{n-j}\binom njj^n=n!$ and that $\sum_{j=0}^n(-1)^j\binom njj^m=0$ for $0\le m<n$. Prove the
second fact again with Theorem~\ref{thm:newton}.
\end{exercise}

\begin{exercise}[$\star$]\label{ex:engine}
A polynomial of degree $3$ takes the values $1,3,11,31,69$ at $x=0,1,2,3,4$. Find it from the difference table, then
check it is $x^3+x+1$ (Figure~\ref{fig:difftable}). Why do four values suffice, and what does the fifth one check?
\end{exercise}

\begin{exercise}[$\star\star$]\label{ex:powersmodp}
Let $p$ be a prime and $m\ge1$. Show that
\[
\sum_{k=1}^{p-1}k^m\equiv\begin{cases}-1\pmod p&\text{if }(p-1)\mid m,\\0\pmod p&\text{otherwise.}\end{cases}
\]
\end{exercise}

\begin{exercise}[$\star\star\star$]\label{ex:staudt}
\emph{The theorem of von Staudt and Clausen (1840).} Let $m\ge2$ be even. This exercise proves that
\[
B_m+\sum_{p\text{ prime},\ (p-1)\mid m}\frac1p\ \in\ \Z.
\]
For example $B_2+\frac12+\frac13=1$ and $B_{12}+\frac12+\frac13+\frac15+\frac17+\frac1{13}=1$.
\begin{enumerate}[label=(\alph*)]
\item Show that $B_m=\sum_{k=0}^m\frac{(-1)^kk!}{k+1}\stirlingii mk$ for $m\ge1$. (Start from
$\frac t{e^t-1}=\frac{\ln(1+u)}u$ with $u=e^t-1$, the formal series $\frac{\ln(1+u)}u=\sum_k\frac{(-1)^ku^k}{k+1}$, and
$\frac{(e^t-1)^k}{k!}=\sum_n\stirlingii nk\frac{t^n}{n!}$.)
\item Show that if $k+1$ is composite and $k+1\ne4$, then $\frac{k!}{k+1}$ is an integer.
\item Show that $\stirlingii m3$ is even for even $m\ge2$, so the term $k=3$ is an integer.
\item For $k+1=p$ prime, show that $(p-1)!\stirlingii m{p-1}\equiv\sum_{j=1}^{p-1}j^m\pmod p$, and conclude with
Exercise~\ref{ex:powersmodp}.
\end{enumerate}
\end{exercise}

\begin{exercise}[$\star\star$]\label{pb:fixmoments}
Let $F$ be the number of fixed points of a uniformly random permutation $\sigma$ of $\{1,\dots,n\}$, and let $\mathfrak S_n$
act on $\{1,\dots,n\}^k$ by $\sigma\cdot(x_1,\dots,x_k)=(\sigma(x_1),\dots,\sigma(x_k))$.
\begin{enumerate}[label=(\alph*)]
\item Show that $\mathbb E[F^k]$ is the number of orbits of this action.
\item Deduce that $\mathbb E[F^k]=\sum_{j=0}^n\stirlingii kj$, and that $\mathbb E[F^k]$ is the Bell number $B_k$ when $k\le n$.
\item Deduce that the first $n$ moments of $F$ are exactly those of a Poisson variable of parameter $1$.
\end{enumerate}
\end{exercise}

\begin{exercise}[$\star\star$]\label{pb:necklacering}
For an integer $a\ge1$, let $M(n,a)$ be the number of primitive necklaces of length $n$ over an alphabet of $a$ letters:
necklaces whose $n$ rotations are pairwise distinct.
\begin{enumerate}[label=(\alph*)]
\item Show that $\sum_{d\mid n}d\,M(d,a)=a^n$. Deduce that $M(n,a)=\frac1n\sum_{d\mid n}\mu(d)a^{n/d}$, and that
$\sum_{d\mid n}\mu(d)a^{n/d}\equiv0\pmod n$ for \emph{every} integer $a$, negative ones included.
\item Prove the \emph{cyclotomic identity}, in $\Z[[x]]$:
\[
\frac1{1-ax}=\prod_{n\ge1}\Bigl(\frac1{1-x^n}\Bigr)^{M(n,a)}.
\]
\item (Metropolis and Rota.) Show that for all $a,b\ge1$,
\[
M(n,ab)=\sum_{\operatorname{lcm}(i,j)=n}\gcd(i,j)\,M(i,a)\,M(j,b),
\]
and explain the identity with necklaces over the alphabet of pairs of letters.
\end{enumerate}
\end{exercise}

\section{Hints}

\begin{hint}{ex:k22k}
Sum by parts with $u=k^2$, $\Delta v=2^k$, then use Example~\ref{ex:k2k}. Mind the range: the exercise sums up to $n$
included.
\end{hint}
\begin{hint}{ex:triple}
The summand is $(k-1)^{\underline{-3}}$. Apply Proposition~\ref{prop:falling} with $m=-2$.
\end{hint}
\begin{hint}{ex:kHk}
$k=k^{\underline1}=\Delta\frac{k^{\underline2}}2$. Sum by parts with $u=H_k$, and use $\frac{(k+1)^{\underline2}}{k+1}=k$.
\end{hint}
\begin{hint}{ex:surjections}
(a) A surjection is a partition of the source into $k$ fibres plus a bijection from the fibres to the target. For the
second count, subtract the maps that miss a given set of $j$ points of the target.
(b) There is no surjection from an $m$-set onto an $n$-set when $m<n$.
\end{hint}
\begin{hint}{ex:engine}
The first column of differences is $2,8,20,38$. A cubic has constant third differences.
\end{hint}
\begin{hint}{ex:powersmodp}
Take a primitive root $g$ modulo $p$ (Volume~2, Chapter~5): the $k^m$ are the $g^{jm}$, a geometric sum.
\end{hint}
\begin{hint}{ex:staudt}
(a) Substitute and compare the coefficients of $\frac{t^m}{m!}$. (b) Write $k+1=ab$ and find $a$ and $b$, or $a$ and $2a$,
among $1,\dots,k$. (c) $6\stirlingii m3=3^m-3\cdot2^m+3$, from the explicit formula of the last section. (d) Expand with the
explicit formula and use $\binom{p-1}j\equiv(-1)^j\pmod p$.
\end{hint}

\begin{hint}{pb:fixmoments}
Burnside's lemma. An orbit of $k$-tuples is determined by which coordinates are equal.
\end{hint}

\begin{hint}{pb:necklacering}
(a) A word of length $n$ has a primitive period $d\mid n$. (b) Take logarithms. (c) Use (a) for $ab$, and $\gcd\cdot\operatorname{lcm}=ij$.
\end{hint}

\section{Solutions}

\begin{solution}{ex:k22k}
\idea One power of $k$ is removed at each summation by parts, and $2^k$ survives every time, since it is its own
difference. Two summations by parts bring $k^2$ down to a constant.
\writeup Let $S=\sum_{k=0}^nk^22^k=\sum_0^{n+1}u\,\Delta v$ with $u(k)=k^2$ and $v(k)=2^k$. Then $\Delta u=2k+1$ and
$Ev(k)=2^{k+1}$, so by Proposition~\ref{prop:byparts}
\[
\begin{aligned}S&=(n+1)^22^{n+1}-\sum_{k=0}^n(2k+1)2^{k+1}\\&=(n+1)^22^{n+1}-4\sum_{k=0}^nk2^k-2\sum_{k=0}^n2^k.\end{aligned}
\]
By Example~\ref{ex:k2k} with $n+1$ in place of $n$, $\sum_{k=0}^nk2^k=(n-1)2^{n+1}+2$, and $\sum_{k=0}^n2^k=2^{n+1}-1$.
Hence
\[
\begin{aligned}S&=2^{n+1}\bigl((n+1)^2-4(n-1)-2\bigr)-8+2\eqbr=(n^2-2n+3)2^{n+1}-6.\end{aligned}
\]
Check: $n=2$ gives $0+2+16=18=3\cdot8-6$.
\end{solution}

\begin{solution}{ex:triple}
\idea The product of three consecutive integers in the denominator is a negative falling power, whose antidifference is
the next one up.
\writeup $\frac1{k(k+1)(k+2)}=(k-1)^{\underline{-3}}$. By Proposition~\ref{prop:falling}, $\Delta x^{\underline{-2}}=-2x^{\underline{-3}}$, so
$x^{\underline{-3}}=\Delta\bigl(-\frac12x^{\underline{-2}}\bigr)$ and
\[
\begin{aligned}\sum_{k=1}^n(k-1)^{\underline{-3}}&=-\frac12\Bigl(n^{\underline{-2}}-0^{\underline{-2}}\Bigr)\eqbr=\frac14-\frac1{2(n+1)(n+2)}.\end{aligned}
\]
As $n\to\infty$ the sum tends to $\frac14$. Check: $n=1$ gives $\frac16=\frac14-\frac1{12}$.
\end{solution}

\begin{solution}{ex:kHk}
\idea $H_k$ is the discrete logarithm and $k$ a power, so this is $\int x\ln x\dd x$ in disguise, done by parts with the
power integrated.
\writeup Take $u(k)=H_k$ and $v(k)=\frac{k^{\underline2}}2$, so $\Delta v=k$, $\Delta u=\frac1{k+1}$ and $Ev(k)=\frac{(k+1)k}2$. Then
\[
\begin{aligned}\sum_{k=0}^{n-1}kH_k&=\frac{n(n-1)}2H_n-\sum_{k=0}^{n-1}\frac k2\eqbr=\frac{n(n-1)}2H_n-\frac{n(n-1)}4,\end{aligned}
\]
which is the formula. Check: for $n=3$, $0+1\cdot1+2\cdot\frac32=4$ and $3\bigl(\frac{11}6-\frac12\bigr)=4$.
\end{solution}

\begin{solution}{ex:surjections}
\idea Count the same set twice; for (b), a polynomial identity checked by counting needs the counting to be empty.
\writeup (a) A surjection $f:N\to K$ with $|N|=n$, $|K|=k$ determines the partition of $N$ into its $k$ non-empty fibres,
and a bijection between the set of fibres and $K$; conversely each such pair gives one surjection. So there are
$k!\stirlingii nk$. For the second count, for $A\subseteq K$ let $M_A$ be the set of maps $N\to K\setminus A$, of size
$(k-|A|)^n$. By inclusion--exclusion (Volume~1, Chapter~10) the maps missing no point number
$\sum_{A}(-1)^{|A|}(k-|A|)^n=\sum_{i=0}^k(-1)^i\binom ki(k-i)^n$; put $j=k-i$.
(b) With $k=n$, the surjections from an $n$-set onto an $n$-set are the $n!$ bijections. With $n$ in place of $k$ and $m$
in place of $n$, the number of surjections from an $m$-set onto an $n$-set is $0$ when $m<n$, so
$\sum_j(-1)^{n-j}\binom njj^m=0$, and multiplying by $(-1)^n$ gives the claim. With Theorem~\ref{thm:newton}: the sum
is $(-1)^n\Delta^n[x^m]_{x=0}$, and $\Delta^n$ kills polynomials of degree $m<n$.
\end{solution}

\begin{solution}{ex:engine}
\idea Differences of a cubic reach a constant at the third step; the left edge of the table is the list of Newton
coefficients.
\writeup The table is $1,3,11,31,69$; then $2,8,20,38$; then $6,12,18$; then $6,6$. The left edge $1,2,6,6$ gives
$f(x)=1+2x+6\binom x2+6\binom x3=1+2x+3x(x-1)+x(x-1)(x-2)=x^3+x+1$. Four values determine a polynomial of degree at most
$3$ (Theorem~\ref{thm:newton} uses only $f(0),\dots,f(3)$), and the fifth checks that the degree is indeed at most~$3$:
it is what makes the two third differences equal.
\end{solution}

\begin{solution}{ex:powersmodp}
\idea Modulo $p$, the non-zero residues are the powers of a primitive root, so the sum is geometric.
\writeup Let $g$ be a primitive root modulo $p$, so that $k=g^j$ for $j=0,\dots,p-2$ runs over the non-zero residues.
Then $\sum_{k=1}^{p-1}k^m\equiv\sum_{j=0}^{p-2}(g^m)^j$. If $(p-1)\mid m$ then $g^m\equiv1$ and the sum is $p-1\equiv-1$.
Otherwise $g^m\not\equiv1$, and $(g^m-1)\sum_{j=0}^{p-2}g^{mj}=g^{m(p-1)}-1\equiv0$ by Fermat's little theorem; since
$g^m-1$ is invertible modulo $p$, the sum is $\equiv0$.
\end{solution}

\begin{solution}{ex:staudt}
\idea Express $B_m$ as a sum of terms $\frac{k!}{k+1}\times$ integer: such a term is an integer unless $k+1$ is prime (or
$4$), and in the prime case its fractional part is decided by sums of powers modulo $p$, which the previous exercise
knows.
\writeup (a) In formal power series, $u=e^t-1$ has zero constant term, so substituting it into
$\frac{\ln(1+u)}u=\sum_{k\ge0}\frac{(-1)^ku^k}{k+1}$ is legitimate (only finitely many $k$ contribute to each power of $t$), and
gives the series of $\frac t{e^t-1}$ since $\ln(1+u)=\ln e^t=t$ (Chapter~2 justifies substitution into formal series
and this identity). By Theorem~\ref{thm:stirling} and the explicit formula,
\[
\begin{aligned}\frac{u^k}{k!}&=\frac1{k!}\sum_{j=0}^k\binom kj(-1)^{k-j}e^{jt}\\&=\sum_n\Bigl(\frac1{k!}\sum_j(-1)^{k-j}\binom kjj^n\Bigr)\frac{t^n}{n!}\\&=\sum_n\stirlingii nk\frac{t^n}{n!}.\end{aligned}
\]
So $\frac t{e^t-1}=\sum_k\frac{(-1)^kk!}{k+1}\sum_n\stirlingii nk\frac{t^n}{n!}$, and the coefficient of $\frac{t^m}{m!}$ is the
formula (terms with $k>m$ vanish).

(b) Let $k+1=ab$ with $2\le a\le b\le k$. If $a<b$, the distinct factors $a$ and $b$ of $k!$ give $ab\mid k!$. If $a=b$,
then $a\ge3$ since $k+1\ne4$, and $a<2a<a^2=k+1$, so $a$ and $2a$ are distinct factors of $k!$, and $a^2\mid2a^2\mid k!$.

(c) The explicit formula gives $6\stirlingii m3=3^m-3\cdot2^m+3$. For even $m\ge2$, $3^m=9^{m/2}\equiv1\pmod 4$ and
$3\cdot2^m\equiv0\pmod 4$, so $6\stirlingii m3\equiv4\equiv0\pmod 4$, hence $\stirlingii m3$ is even. The term $k=3$ is
$-\frac64\stirlingii m3=-3\cdot\frac{\stirlingii m3}2\in\Z$.

(d) For $k+1=p$ prime,
\[
(p-1)!\stirlingii m{p-1}=\sum_{j=0}^{p-1}(-1)^{p-1-j}\binom{p-1}jj^m.
\]
In $\binom{p-1}j=\prod_{i=1}^j\frac{p-i}i$
each factor is $\equiv-1$ modulo $p$, so $\binom{p-1}j\equiv(-1)^j$, and $(-1)^{p-1-j}(-1)^j=(-1)^{p-1}\equiv1\pmod p$
(for $p=2$, $-1\equiv1$). The term $j=0$ vanishes since $m\ge1$. So $(p-1)!\stirlingii m{p-1}\equiv\sum_{j=1}^{p-1}j^m$,
which by Exercise~\ref{ex:powersmodp} is $\equiv-1$ if $(p-1)\mid m$ and $\equiv0$ otherwise. Hence the term of (a)
with $k=p-1$, namely $\frac{(-1)^{p-1}(p-1)!\stirlingii m{p-1}}p$, is an integer plus $-\frac1p$ when $(p-1)\mid m$
(for odd $p$ the sign $(-1)^{p-1}$ is $1$; for $p=2$ it does not matter, since $\frac12\equiv-\frac12$ modulo $\Z$), and an
integer otherwise.

Putting (a)--(d) together: the terms with $k+1$ composite, including $k+1=4$, and the term $k=0$ (which is $0$ for
$m\ge1$) are integers, and each prime $p$ with $(p-1)\mid m$ contributes $-\frac1p$ modulo an integer. So
$B_m+\sum_{(p-1)\mid m}\frac1p\in\Z$. (A prime $p$ with $p-1>m$ gives no term in (a), since $k\le m$ there, and it is
not in the sum either, since $(p-1)\nmid m$.)
\end{solution}

\begin{solution}{pb:fixmoments}
\idea $F(\sigma)^k$ counts the $k$-tuples of fixed points of $\sigma$, that is the $k$-tuples fixed by $\sigma$ acting on
$\{1,\dots,n\}^k$. Averaging a number of fixed points over a group is exactly what Burnside's lemma turns into a number of
orbits. The orbits are equality patterns, that is set partitions, and the Bell numbers are the moments of the Poisson law.

\writeup (a) A tuple $x$ is fixed by $\sigma$ iff $\sigma(x_i)=x_i$ for every $i$, iff every $x_i$ is a fixed point of
$\sigma$; there are $F(\sigma)^k$ such tuples. By Burnside's lemma (Volume~3, Chapter~3), the number of orbits is
$\frac1{n!}\sum_\sigma F(\sigma)^k=\mathbb E[F^k]$.

(b) To a tuple $x$ attach the partition $\pi(x)$ of $\{1,\dots,k\}$ whose blocks are the sets $\{i:x_i=v\}$, $v$ a value taken
by $x$. A permutation does not change it, since $\sigma(x_i)=\sigma(x_j)$ iff $x_i=x_j$. Conversely, if $\pi(x)=\pi(y)$, the map
$x_i\mapsto y_i$ is a well-defined bijection from the set of values of $x$ onto that of $y$, and extends to a permutation
$\sigma$ of $\{1,\dots,n\}$ (both complements have the same size), with $\sigma\cdot x=y$. So the orbits correspond to the
partitions of $\{1,\dots,k\}$ into at most $n$ blocks (a tuple takes at most $n$ values, and every such partition occurs).
There are $\stirlingii kj$ partitions into $j$ blocks (Volume~5, Chapter~1), so $\mathbb E[F^k]=\sum_{j\le n}\stirlingii kj$,
which is $\sum_j\stirlingii kj=B_k$ when $k\le n$, since then no partition has more than $n$ blocks.

(c) If $N$ is Poisson of parameter $1$, $\mathbb E[N^k]=e^{-1}\sum_{m\ge0}\frac{m^k}{m!}=B_k$ by Dobiński's formula (Volume~5,
Chapter~2). So $\mathbb E[F^k]=\mathbb E[N^k]$ for $1\le k\le n$; for $k=n+1$ they differ, since $\mathbb E[F^{n+1}]=B_{n+1}-1$
(the partition into $n+1$ singletons is missing).
\tools Volume~3, Chapter~3 (Burnside's lemma); Volume~5, Chapters~1--2 (Stirling and Bell numbers, Dobiński); Volume~4, Chapter~2.
\leads Moments that agree up to order $n$ imply a total variation distance between $F$ and Poisson(1) of order
$\frac{2^n}{n!}$, super-exponentially small (Diaconis and Pitman; Arratia and Tavaré, 1992). A numerical calculation compares the two laws.
\end{solution}

\begin{solution}{pb:necklacering}
\idea Every word of length $n$ is a primitive word of some length $d\mid n$ repeated $\frac nd$ times, and a primitive necklace
of length $d$ contains exactly $d$ words. That gives (a), and everything else is (a) in disguise: the logarithm of the product
in (b), and the factorisation $(ab)^n=a^nb^n$ in (c).

\writeup (a) A word $w$ of length $n$ has a smallest period $d$, the least $d\ge1$ such that rotating by $d$ fixes $w$. The
rotations that fix $w$ form a subgroup of $\Z/n\Z$, generated by $d$, so $d\mid n$, and $w$ is $u^{n/d}$ with $u$ a word of length
$d$ whose $d$ rotations are distinct. Conversely each such $u$ gives a word of smallest period $d$. The words $u$ of length $d$
with distinct rotations fall into primitive necklaces of exactly $d$ words. So $a^n=\sum_{d\mid n}d\,M(d,a)$, and Möbius
inversion (Volume~2, Chapter~4) gives $nM(n,a)=\sum_{d\mid n}\mu(d)a^{n/d}$.

The polynomial $P(t)=\frac1n\sum_{d\mid n}\mu(d)t^{n/d}$ has rational coefficients and takes integer values at all $t\ge1$, and also $P(0)=0$.
A polynomial with this property is an integer combination of the binomial polynomials $\binom tj$ (Newton's interpolation
formula, Volume~5, Chapter~1), and these take integer values at negative integers too. So $\sum_{d\mid n}\mu(d)a^{n/d}\equiv0\pmod n$
for every $a\in\Z$.

(b) Both sides have constant term $1$, so it is enough to compare logarithms (Volume~5, Chapter~2):
\[
\log\prod_{n\ge1}(1-x^n)^{-M(n,a)}=\sum_{n\ge1}M(n,a)\sum_{m\ge1}\frac{x^{nm}}m=\sum_{N\ge1}\frac{x^N}N\sum_{n\mid N}n\,M(n,a)
=\sum_{N\ge1}\frac{(ax)^N}N,
\]
which is $-\log(1-ax)$, by (a) with $N=nm$.

(c) Write $R(n)$ for the right-hand side. For $n\ge1$,
\[
\sum_{d\mid n}d\,R(d)=\sum_{d\mid n}\ \sum_{\operatorname{lcm}(i,j)=d}\operatorname{lcm}(i,j)\gcd(i,j)M(i,a)M(j,b)=\sum_{i\mid n,\ j\mid n}ij\,M(i,a)M(j,b),
\]
because $\operatorname{lcm}(i,j)\mid n$ if and only if $i\mid n$ and $j\mid n$, and $\operatorname{lcm}\cdot\gcd=ij$. By (a) this is
$a^nb^n=(ab)^n=\sum_{d\mid n}d\,M(d,ab)$. Möbius inversion of $d\,R(d)$ and $d\,M(d,ab)$ gives $R(n)=M(n,ab)$.

The meaning: a word over the alphabet of pairs is a pair of words, one over each alphabet. If the first has smallest period $i$
and the second $j$, the pair has smallest period $\operatorname{lcm}(i,j)$. A pair of primitive necklaces of lengths $i$ and $j$
gives $ij$ pairs of words, hence $\frac{ij}{\operatorname{lcm}(i,j)}=\gcd(i,j)$ primitive necklaces of length $\operatorname{lcm}(i,j)$.
\tools Volume~2, Chapter~4 (Möbius inversion); Volume~5, Chapters~1 and~2 (Newton series, logarithms of power series).
\leads $M(n,q)$ is also the number of monic irreducible polynomials of degree $n$ over $\mathbb F_q$ (Volume~3, Chapter~9), and then
(b) is unique factorisation in $\mathbb F_q[X]$. Identity (c) makes the sequences $(M(n,a))_n$ into a ring, the \emph{necklace
ring} of Metropolis and Rota, isomorphic to the ring of Witt vectors that controls $p$-adic arithmetic.
\end{solution}

\section*{Sources and further reading}

\begin{itemize}
\item R.~Graham, D.~Knuth and O.~Patashnik, \emph{Concrete Mathematics}, 2nd ed., Addison-Wesley, 1994. Chapter~2
(finite calculus, summation by parts) and Chapter~6 (Stirling and Bernoulli numbers). The notation
$\stirlingii nk$, $\stirlingi nk$ and $x^{\underline m}$ is theirs.
\item D.~Knuth, \emph{Johann Faulhaber and sums of powers}, Mathematics of Computation 61 (1993), 277--294.
\item J.~Bernoulli, \emph{Ars Conjectandi}, 1713; English translation by E.~D.~Sylla, \emph{The Art of Conjecturing},
Johns Hopkins University Press, 2006. The ``half of a quarter of an hour'' is in Part~II, Chapter~3.
\item G.~Boole, \emph{A Treatise on the Calculus of Finite Differences}, Macmillan, 1860; C.~Jordan, \emph{Calculus of
Finite Differences}, 1939, reprinted by Chelsea, 1965.
\item K.~Ireland and M.~Rosen, \emph{A Classical Introduction to Modern Number Theory}, 2nd ed., Springer, 1990,
Chapter~15 (Bernoulli numbers, von Staudt--Clausen, Kummer and irregular primes).
\item J.~H.~Conway and R.~K.~Guy, \emph{The Book of Numbers}, Copernicus, 1996, Chapter~4 (Bernoulli numbers and
figurate numbers).
\item D.~Swade, \emph{The Difference Engine: Charles Babbage and the Quest to Build the First Computer}, Viking, 2001.
\item J.~Bernoulli, \emph{Ars Conjectandi}, Basel, 1713, Part~2, Chapter~3; English translation by E.~D.~Sylla, \emph{The Art of Conjecturing}, Johns Hopkins University Press, 2006, for the quotation.
\end{itemize}
